% =====================================================================
\appendix
\section{Certification table and comparison data}
\label{app:tables}
% =====================================================================

\subsection{The certification table}

Table~\ref{tab:cert} is the complete output of the certification stage: all
$28$ certified codes, reproduced verbatim from the computation log
(\filepath{04_COMPUTATION/phase7_extended_search_log.txt}, Part~6). The
columns M1, M2, M3 record the three independent distance methods of
Section~\ref{sec:cert}. At most two of the three are run for any one code:
a dash in the M1 column means M1 could not complete within its enumeration
cap, in which case the support search M2 is run instead; ``n/a'' in the M2
column means M2 was not run, because M1 had already finished and M2 is run
only when M1 cannot. Thus exactly one of M1 and M2 carries the exact distance
for each code. A ``$>2$'' entry means the ambient quantum search
(M3) reached weight $2$ without finding a vector, i.e. it certifies $d\ge3$ as
a \emph{lower} bound only; M3 is therefore a cross-check rather than an exact
method. The column ``exact'' says yes when every method that completed agrees
with the reported value, the stabilizer is $\F_p$-abelian, and $K$ and $c$
recompute from $k$ and $h_E$.

\begin{table}[ht]
\centering
\small
\caption{All $28$ certified codes (Part~6 of the Phase-7 log, verbatim).
$S$-def is the entanglement-assisted Singleton defect $N-K+c-2(d-1)$.}
\label{tab:cert}
\begin{tabular}{@{}llcrrrrrl@{}}
\toprule
code & family & route & $k$ & $K$ & $d$ & $c$
 & $S$-def & M1 / M2 / M3\\
\midrule
$[[24,8,4]]_9$   & $(9,1,6)$ & SO & 8  & 8  & 4  & 0  & 10 & -- / 4 / $>2$\\
$[[20,4,4]]_4$   & $(4,1,5)$ & SO & 8  & 4  & 4  & 0  & 10 & -- / 4 / $>2$\\
$[[20,4,4]]_4$   & $(4,1,5)$ & DC & 12 & 4  & 4  & 0  & 10 & -- / 4 / $>2$\\
$[[20,2,10;18]]_4$ & $(4,1,5)$ & EA & 18 & 2  & 10 & 18 & 18 & 10 / n/a / $>2$\\
$[[24,2,10;18]]_9$ & $(9,1,6)$ & EA & 20 & 2  & 10 & 18 & 22 & 10 / n/a / $>2$\\
$[[4,2,2;2]]_4$  & $(4,1,1)$ & EA & 2  & 2  & 2  & 2  & 2  & 2 / n/a / 2\\
$[[8,4,3;4]]_4$  & $(4,1,2)$ & EA & 4  & 4  & 3  & 4  & 4  & 3 / n/a / 3\\
$[[12,2,6;10]]_4$ & $(4,1,3)$ & EA & 10 & 2  & 6  & 10 & 10 & 6 / n/a / $>2$\\
$[[16,4,5;8]]_4$ & $(4,1,4)$ & EA & 10 & 4  & 5  & 8  & 12 & 5 / n/a / $>2$\\
$[[24,4,9;20]]_4$ & $(4,1,6)$ & EA & 20 & 4  & 9  & 20 & 24 & 9 / n/a / $>2$\\
$[[8,4,2;4]]_4$  & $(4,2,1)$ & EA & 4  & 4  & 2  & 4  & 6  & 2 / n/a / 2\\
$[[16,4,4;12]]_4$ & $(4,2,2)$ & EA & 12 & 4  & 4  & 12 & 18 & 4 / n/a / $>2$\\
$[[24,4,6;20]]_4$ & $(4,2,3)$ & EA & 20 & 4  & 6  & 20 & 30 & 6 / n/a / $>2$\\
$[[12,6,2;6]]_4$ & $(4,3,1)$ & EA & 6  & 6  & 2  & 6  & 10 & 2 / n/a / 2\\
$[[16,6,2;10]]_4$ & $(4,4,1)$ & EA & 10 & 6  & 2  & 10 & 18 & 2 / n/a / 2\\
$[[4,2,2;2]]_8$  & $(8,1,1)$ & EA & 2  & 2  & 2  & 2  & 2  & 2 / n/a / 2\\
$[[8,4,3;4]]_8$  & $(8,1,2)$ & EA & 4  & 4  & 3  & 4  & 4  & 3 / n/a / 3\\
$[[12,2,6;10]]_8$ & $(8,1,3)$ & EA & 10 & 2  & 6  & 10 & 10 & 6 / n/a / $>2$\\
$[[12,4,2;8]]_8$ & $(8,3,1)$ & EA & 8  & 4  & 2  & 8  & 14 & 2 / n/a / 2\\
$[[4,2,2;2]]_9$  & $(9,1,1)$ & EA & 2  & 2  & 2  & 2  & 2  & 2 / n/a / 2\\
$[[8,2,4;6]]_9$  & $(9,1,2)$ & EA & 6  & 2  & 4  & 6  & 6  & 4 / n/a / $>3$\\
$[[12,2,5;6]]_9$ & $(9,1,3)$ & EA & 8  & 2  & 5  & 6  & 8  & 5 / n/a / $>2$\\
$[[16,2,8;14]]_9$ & $(9,1,4)$ & EA & 14 & 2  & 8  & 14 & 14 & 8 / n/a / $>2$\\
$[[20,4,9;16]]_9$ & $(9,1,5)$ & EA & 16 & 4  & 9  & 16 & 16 & 9 / n/a / $>2$\\
$[[8,4,2;4]]_9$  & $(9,2,1)$ & EA & 4  & 4  & 2  & 4  & 6  & 2 / n/a / 2\\
$[[16,4,4;12]]_9$ & $(9,2,2)$ & EA & 12 & 4  & 4  & 12 & 18 & 4 / n/a / $>2$\\
$[[12,4,2;8]]_9$ & $(9,3,1)$ & EA & 8  & 4  & 2  & 8  & 14 & 2 / n/a / 2\\
$[[24,6,4;18]]_9$ & $(9,3,2)$ & EA & 18 & 6  & 4  & 18 & 30 & -- / 4 / $>2$\\
\bottomrule
\end{tabular}
\end{table}

\subsection{The non-principal witnesses}

The witnesses to Theorem~\ref{thm:nonprincipal} and the quantum codes they
produce, from Part~5 of the same log:

\begin{center}
\begin{tabular}{@{}ll@{}}
\toprule
witness & codes produced\\\midrule
$(9,2,3)$, $s=0$ & R1: $[[24,14,2]]_9$ (abelian); R3: $[[24,6,2;2]]_9$\\
$(9,2,3)$, $s=1$ & R1: $[[24,14,2]]_9$ (abelian); R3: $[[24,6,2;2]]_9$\\
$(4,2,4)$, $s=0$ & R3: $[[32,8,2;4]]_4$\\
$(4,2,4)$, $s=1$ & R3: $[[32,8,2;4]]_4$\\
\bottomrule
\end{tabular}
\end{center}

In all four cases $\MU_{\Bq}(W)=3$ and $\MU_{\Bq}(\Hull_s(W))=3$. The two
$(9,2,3)$ witnesses additionally have $\Hull_s(W)=W$, i.e. the code equals its
own Galois hull.

\subsection{Best known distances used in Section~\ref{sec:cmp}}

Table~\ref{tab:grassl} is the raw data retrieved from
\cite{GrasslCodetables} on 2026-11-04, in the form lower..upper bound for
$[[N,K]]_q$. The entry ``not available'' is the site's own response for
$q=8$.

\begin{table}[ht]
\centering
\caption{Best known stabilizer distances, codetables.de, retrieved
2026-11-04. The site prints the \emph{prime} $p$ in its code labels, so
$q=4$ appears as $[[N,K]]_2$ and $q=9$ as $[[N,K]]_3$.}
\label{tab:grassl}
\begin{tabular}{@{}cc@{\qquad}cc@{}}\toprule
$[[N,K]]_q$ & lower..upper & $[[N,K]]_q$ & lower..upper\\\midrule
$[[16,4]]_4$   & 5..5  & $[[8,4]]_9$   & 3..3\\
$[[16,12]]_4$  & 2..2  & $[[16,8]]_9$  & 4..4\\
$[[20,4]]_4$   & 6..6  & $[[24,8]]_9$  & 6..8\\
$[[20,12]]_4$  & 3..3  & $[[24,10]]_9$ & 6..7\\
$[[20,16]]_4$  & 2..2  & $[[24,12]]_9$ & 5..6\\
$[[24,16]]_4$  & 3..3  & $[[8,4]]_8$   & not available\\
$[[24,20]]_4$  & 2..2  & $[[12,2]]_8$  & not available\\
\bottomrule
\end{tabular}
\end{table}
